\section{Desafios Críticos e Impacto Social}
\label{sec:ia-desafios}

Substituir ontologias especulativas e fluxos mercadológicos entusiasmados pela dura concretude da clínica em larga escala, e especialmente dentro do escopo logístico do Sistema Único de Saúde (SUS), demanda a resolutividade de desafios técnicos severos, desigualdades econômicas latentes e barreiras regulatórias imponentes \cite{frota2025, du2020}.

\textbf{Escassez de Generalização e Dados Sujos:} Arquiteturas treinadas em recortes amostrais oriundos de laboratórios endinheirados do hemisfério norte frequentemente colapsam sob o peso da realidade de equipamentos de menor resolução técnica comumente encontrados na periferia clínica. A validação multicêntrica transcontinental e rigorosa se apresenta como o antídoto único ao \textit{overfitting} e ao colapso de predição inter-populacional \cite{khoshfekr2025}.

\textbf{Exclusão e Desertos Digitais:} A concentração endêmica de especialistas nos grandes centros urbanos já caracteriza um profundo gargalo no cuidado oral. A mera digitalização sem responsabilidade social tende a aprofundar essa ruptura sistêmica. Gêmeos digitais não podem permanecer como artigos de elite acessíveis apenas à parcela abastada da população. Como discutido no trabalho de Cuadros et al. (2023), comunidades vulnerabilizadas operam sob uma gravosa falha estrutural de infraestrutura digital \cite{cuadros2023}. A solução jaz em arquiteturas e protocolos orquestrados \textit{open-source} em nuvens otimizadas, e na expansão em larga escala da infraestrutura de IoT tele-odontológica.

\destaque{
A compressão inteligente de modelos e a otimização algorítmica — conduzidas por subsistemas como o "Structural Noise Scanner" de plataformas em nuvem — mostram-se decisivas para operar grandes volumes de dados (como os DTs) através de conexões restritas e maquinário de menor latência presente nos postos de saúde mais remotos.
}

\textbf{Aparato Regulatório:} No escopo da vigilância sanitária moderna, uma \textit{software as a medical device} (SaMD) demanda validação ininterrupta. A morosidade da atualização da regulação sanitária no Brasil, que hoje delineia seus primeiros passos formais via ISO/IEC 3173 \cite{khoshfekr2025}, restringe os impulsos de pesquisadores e sufoca empresas emergentes por uma aguda falta de lastro metodológico sancionado. O ecossistema ideal será aquele dotado de certificação clínica e protocolos automatizados de validação de qualidade e rastreabilidade algorítmica por \textit{design} \cite{aisociety2026}.
